\documentclass[12pt, a4paper]{article}

% ===== 基础宏包 =====
\usepackage[UTF8]{ctex}          % 中文支持
\usepackage[
  top=2.5cm, bottom=2.5cm,
  left=3cm,  right=3cm
]{geometry}
\usepackage{amsmath}
\usepackage{amssymb}

% ===== 代码高亮 =====
\usepackage{listings}
\usepackage{xcolor}

\definecolor{codebg}{RGB}{245,245,245}
\definecolor{codecomment}{RGB}{100,100,100}
\definecolor{codekeyword}{RGB}{0,0,180}
\definecolor{codestring}{RGB}{163,21,21}
\definecolor{codenumber}{RGB}{128,128,128}
\definecolor{warningbg}{RGB}{255,248,220}
\definecolor{warningborder}{RGB}{220,180,0}
\definecolor{tipbg}{RGB}{235,245,255}
\definecolor{tipborder}{RGB}{70,130,200}

\lstdefinestyle{cpp}{
  language=C++,
  backgroundcolor=\color{codebg},
  basicstyle=\ttfamily\small,
  keywordstyle=\color{codekeyword}\bfseries,
  commentstyle=\color{codecomment}\itshape,
  stringstyle=\color{codestring},
  numberstyle=\tiny\color{codenumber},
  numbers=left,
  numbersep=8pt,
  stepnumber=1,
  breaklines=true,
  breakatwhitespace=false,
  tabsize=4,
  showspaces=false,
  showstringspaces=false,
  frame=single,
  framesep=4pt,
  rulecolor=\color{gray!40},
  xleftmargin=1.5em,
  framexleftmargin=1.5em,
  morekeywords={string, stringstream, getline, ios, numeric\_limits, streamsize},
}

\lstdefinestyle{input}{
  backgroundcolor=\color{gray!10},
  basicstyle=\ttfamily\small,
  frame=single,
  rulecolor=\color{gray!40},
  framesep=4pt,
  numbers=none,
  breaklines=true,
  xleftmargin=0.5em,
  framexleftmargin=0.5em,
}

% ===== 警告 / 提示框 =====
\usepackage{tcolorbox}
\tcbuselibrary{skins, breakable}

\newtcolorbox{warningbox}{
  enhanced,
  colback=warningbg,
  colframe=warningborder,
  leftrule=4pt, toprule=0.5pt, bottomrule=0.5pt, rightrule=0.5pt,
  arc=2pt,
  fonttitle=\bfseries,
  title={⚠\quad 经典陷阱},
  breakable,
}

\newtcolorbox{tipbox}{
  enhanced,
  colback=tipbg,
  colframe=tipborder,
  leftrule=4pt, toprule=0.5pt, bottomrule=0.5pt, rightrule=0.5pt,
  arc=2pt,
  fonttitle=\bfseries,
  title={💡\quad 通用加速技巧（必加）},
  breakable,
}

% ===== 表格 =====
\usepackage{booktabs}
\usepackage{array}
\usepackage{tabularx}

% ===== 超链接（目录可点击）=====
\usepackage{hyperref}
\hypersetup{
  colorlinks=true,
  linkcolor=blue!60!black,
  urlcolor=blue!60!black,
}

% ===== 页眉页脚 =====
\usepackage{fancyhdr}
\pagestyle{fancy}
\fancyhf{}
\rhead{\small OJ 输入模式速查}
\lhead{\small C++ 算法竞赛}
\cfoot{\thepage}
\renewcommand{\headrulewidth}{0.4pt}

% ===== 章节格式 =====
\usepackage{titlesec}
\titleformat{\section}[block]
  {\large\bfseries}
  {\thesection}{0.8em}{}
  [\vspace{-0.3em}\rule{\linewidth}{0.6pt}]

\titleformat{\subsection}[block]
  {\normalsize\bfseries}
  {\thesubsection}{0.6em}{}

% ===== 其他 =====
\usepackage{parskip}
\setlength{\parskip}{0.5em}
\usepackage{enumitem}
\setlist{noitemsep, topsep=4pt}

% ========================================
\begin{document}

\begin{titlepage}
  \centering
  \vspace*{3cm}
  {\Huge\bfseries OJ 输入模式速查手册}\par
  \vspace{0.8cm}
  {\large C++ 算法竞赛 · 国内大学课程 OJ 适用}\par
  \vspace{0.4cm}
  {\normalsize 覆盖 PTA、洛谷、牛客、HDU 等主流平台}\par
  \vspace{2cm}
  \rule{0.6\linewidth}{0.5pt}\par
  \vspace{0.4cm}
  {\small 掌握本文六种模式，可应对 99\% 的 OJ 输入场景}\par
  \vfill
\end{titlepage}

\tableofcontents
\newpage

% ========================================
\section{概述}

在算法竞赛与课程 OJ 中，输入处理往往是解题的第一道门槛。
本文总结了六种最常见的输入模式及其标准解法，并给出了易错点说明。

\textbf{涉及头文件一览：}

\begin{lstlisting}[style=cpp, numbers=none]
#include <iostream>   // cin, cout, getline
#include <string>     // string
#include <sstream>    // stringstream（案例五专用）
#include <limits>     // numeric_limits（cin.ignore 增强版）
using namespace std;
\end{lstlisting}

% ========================================
\section{案例一：不定数量整数输入（EOF 结束）}

\subsection{特征}

题目说明"输入包含多组数据"或"处理到文件结束"，
没有给出数据组数，直到读到文件末尾（EOF）为止。

\subsection{输入示例}

\begin{lstlisting}[style=input]
10 20
30 40
...（直到文件结束）
\end{lstlisting}

\subsection{解决方法}

利用 \texttt{cin >> num} 的返回值——它返回流对象本身；
当遇到 EOF 或读取失败时，流状态变为 \texttt{false}，\texttt{while} 循环自动结束。

\begin{lstlisting}[style=cpp]
#include <iostream>
using namespace std;

int main() {
    int a, b;
    // cin >> a >> b 返回流对象，EOF 时转为 false，循环结束
    while (cin >> a >> b) {
        // 处理 a 和 b
    }
    return 0;
}
\end{lstlisting}

\textbf{扩展：} 每行整数数量不固定时，改为：

\begin{lstlisting}[style=cpp]
int x;
while (cin >> x) {
    // 不管每行几个数，全部逐个读入
}
\end{lstlisting}

% ========================================
\section{案例二：指定组数输入}

\subsection{特征}

第一行是整数 $N$，表示接下来有 $N$ 组数据。

\subsection{输入示例}

\begin{lstlisting}[style=input]
2
10 20
30 40
\end{lstlisting}

\subsection{解决方法}

先读入 $N$，再循环 $N$ 次。

\begin{lstlisting}[style=cpp]
#include <iostream>
using namespace std;

int main() {
    int n;
    cin >> n;
    while (n--) {          // 等价于 for (int i = 0; i < n; i++)
        int a, b;
        cin >> a >> b;
        // 处理数据
    }
    return 0;
}
\end{lstlisting}

% ========================================
\section{案例三：多组数据 + 每组内含 N 个数}

\subsection{特征}

这是国内课程 OJ（PTA、洛谷入门）中极为高频的嵌套结构：
第一行为组数 $T$，每组第一行是本组数据量 $N$，
随后 $N$ 行（或 $N$ 个数）是具体数据。

\subsection{输入示例}

\begin{lstlisting}[style=input]
2
3
1 2 3
4
10 20 30 40
\end{lstlisting}

\subsection{解决方法}

\begin{lstlisting}[style=cpp]
#include <iostream>
using namespace std;

int main() {
    int t;
    cin >> t;
    while (t--) {
        int n;
        cin >> n;
        for (int i = 0; i < n; i++) {
            int x;
            cin >> x;
            // 处理第 i 个数据
        }
        // 每组处理完毕
    }
    return 0;
}
\end{lstlisting}

% ========================================
\section{案例四：带空格的字符串（整行读取）}

\subsection{特征}

输入包含空格，需要把整行作为一个整体读入，
例如 \texttt{Hello World}、\texttt{I am a student} 等。

\subsection{输入示例}

\begin{lstlisting}[style=input]
Hello World
I am a student
\end{lstlisting}

\subsection{解决方法}

使用 \texttt{getline}。

\begin{lstlisting}[style=cpp]
#include <iostream>
#include <string>
using namespace std;

int main() {
    string s;
    while (getline(cin, s)) {
        // 处理整行字符串 s
    }
    return 0;
}
\end{lstlisting}

\begin{warningbox}
\textbf{经典陷阱：\texttt{cin} 与 \texttt{getline} 混用}

若题目结构为：第一行整数 $N$，之后 $N$ 行带空格的字符串：

\begin{lstlisting}[style=input]
2
Hello World
I love OJ
\end{lstlisting}

\texttt{cin >> n} 读取整数后，换行符 \texttt{'\textbackslash n'} 仍残留在缓冲区，
下一次 \texttt{getline} 会立即读到一个空行。\textbf{必须先吃掉残留换行符}：

\begin{lstlisting}[style=cpp]
int n;
cin >> n;
cin.ignore();   // 吃掉换行符（简洁版，适合大多数情况）
// 更健壮的写法：
// cin.ignore(numeric_limits<streamsize>::max(), '\n');

string s;
for (int i = 0; i < n; i++) {
    getline(cin, s);   // 此时才能正确读取
}
\end{lstlisting}

\texttt{numeric\_limits<streamsize>::max()} 需引入头文件 \texttt{<limits>}，
可忽略缓冲区中任意多个无关字符，比 \texttt{cin.ignore()} 更稳健。
\end{warningbox}

% ========================================
\section{案例五：逐行读取数量不定的整数}

\subsection{特征}

每行有若干整数，数量未知，需要逐行处理。

\subsection{输入示例}

\begin{lstlisting}[style=input]
1 2 3
4 5
6 7 8 9
\end{lstlisting}

\subsection{解决方法}

\texttt{getline} + \texttt{stringstream}（最推荐的标准写法）。

\begin{lstlisting}[style=cpp]
#include <iostream>
#include <string>
#include <sstream>
using namespace std;

int main() {
    string line;
    while (getline(cin, line)) {   // 1. 逐行读取整行
        stringstream ss(line);     // 2. 将字符串包装为流
        int num;
        while (ss >> num) {        // 3. 像用 cin 一样逐个提取
            // 处理 num
        }
        // 到这里，当前行处理完毕，自动进入下一行
    }
    return 0;
}
\end{lstlisting}

\textbf{原理：} \texttt{stringstream} 将字符串模拟为输入流，
可以用 \texttt{>>} 运算符从中提取任意类型的数据，
读完后流状态变为 \texttt{false}，内层 \texttt{while} 自动退出。

% ========================================
\section{案例六：特定结束标记}

\subsection{特征}

题目说明"当输入为 \texttt{0 0} 时结束"或类似的哨兵值。

\subsection{输入示例}

\begin{lstlisting}[style=input]
1 2
10 20
0 0
\end{lstlisting}

\subsection{解决方法}

在 \texttt{while} 循环内增加判断条件。

\begin{lstlisting}[style=cpp]
#include <iostream>
using namespace std;

int main() {
    int a, b;
    while (cin >> a >> b) {
        if (a == 0 && b == 0) break;   // 遇到哨兵值，跳出
        // 处理数据
    }
    return 0;
}
\end{lstlisting}

\textbf{变体：} 结束标记也可能是单个值（如 $-1$）或字符串（如 \texttt{END}），
原理相同，修改判断条件即可。

% ========================================
\section{通用加速技巧}

\begin{tipbox}
在数据量较大时，C++ 的 \texttt{cin}/\texttt{cout} 默认需要与 C 语言的
\texttt{scanf}/\texttt{printf} 保持同步，速度较慢。
在算法竞赛中，建议在 \texttt{main} 函数开头加上以下两行（第三行可选）：

\begin{lstlisting}[style=cpp, numbers=none]
ios::sync_with_stdio(false);   // 关闭 C/C++ I/O 同步，cin 加速
cin.tie(0);                    // 解除 cin 与 cout 的绑定，避免自动刷新
// cout.tie(0);                // 几乎无实际效果，可省略
\end{lstlisting}

\textbf{注意：} 开启加速后，\textbf{禁止混用} \texttt{scanf}/\texttt{printf} 与
\texttt{cin}/\texttt{cout}，否则输入输出顺序可能错乱。
\end{tipbox}

% ========================================
\section{总结对照表}

\begin{table}[h!]
\centering
\renewcommand{\arraystretch}{1.4}
\begin{tabularx}{\textwidth}{
  >{\bfseries}l
  l
  X
}
\toprule
输入类型 & 核心代码 & 关键点 \\
\midrule
EOF 结束
  & \texttt{while(cin >> a)}
  & 流状态自动判断，最简洁 \\

指定组数 $N$
  & \texttt{cin >> n; while(n--)}
  & 先读 $N$，再循环 \\

嵌套：$T$ 组 + 每组 $N$ 个
  & 双层循环
  & 国内 OJ 高频，先读 $T$ 再读 $N$ \\

带空格整行
  & \texttt{getline(cin, s)}
  & 与 \texttt{cin} 混用时需 \texttt{cin.ignore()} \\

每行数量不定
  & \texttt{getline} + \texttt{stringstream}
  & 需引入 \texttt{<sstream>} \\

特定结束标记
  & \texttt{while(cin >> a)\{if(...) break;\}}
  & 循环内判断哨兵值 \\
\bottomrule
\end{tabularx}
\caption{六种 OJ 输入模式速查}
\end{table}

\vspace{1em}
\noindent\rule{\linewidth}{0.4pt}
\begin{center}
{\small 掌握以上六种模式，可应对国内大学课程 OJ 中 99\% 的输入场景。}
\end{center}

\end{document}
